\section{Case study}
\label{sec:case_study}
To test the proposed methodology, this study considers the NREL-5MW \cite{NREL_5MW} Horizontal-Axis Wind Turbine (HAWT). As cited before, the rotor is modeled with the ALM to reduce the computational cost of the simulations. Number of nodes are changed according to the mesh refinement, to ensure always that $\epsilon \le 2\Delta x$ and $\frac{\Delta b}{\Delta x} \le 0.75$ \cite{inproceedings}. Selected combinations of $\Delta b$ are reported in Tab.~\ref{tab:ALM_parameters} as a function of the mesh refinement.
\begin{table}[h]
    \centering
    \begin{tabular}{c|c|c}
        \textbf{Mesh $\Delta x$ [m]} & $\Delta b$ [m] & $\frac{\Delta b}{\Delta x}$ \\
        \hline
        $\sim 3.94$ & 32 & $\sim0.5$ \\
        $\sim 2.82$ & 44 & $\sim0.5$ \\    
        $\sim 2.25$ & 56 & $\sim0.5$ \\
        $\sim 1.78$ & 71 & $\sim0.5$ \\
    \end{tabular}
    \caption{ALM parameters for the different mesh refinements.}
    \label{tab:ALM_parameters}
\end{table}
An inflow velocity $U_{ref}$ at hub height (i.e., $90$~m from the ground) of $11\,\mathrm{m/s}$ is taken as the single reference condition on which the closure is trained, while two additional inflow velocities of $8$ and $15\,\mathrm{m/s}$ are simulated solely to test the trained model on operating conditions never seen during training. Consistently with the design premise of the methodology, no data from these two conditions, nor from the coarser grids used for testing, enter the training set at any stage: together they constitute a strict out-of-sample assessment of the learned closure. For each condition the ABL logarithmic profile is evaluated from the corresponding hub-height velocity as reported in Eq.~\ref{eq:ABL_log_equation}:
\begin{equation}
    U(z) = U_{ref} \frac{\ln(\frac{z+z_0}{z_0})}{\ln(\frac{z_{ref}+z_0}{z_0})}
    \label{eq:ABL_log_equation}
\end{equation}

The data generated with the CFD campaign was evaluated against literature results \cite{alessio5-mw, Troldborg} as also done in previous studies (e.g., \cite{Carloni2026}). Hence, the results are not reported here for brevity, but the interested reader can refer to the cited works for a detailed validation of the CFD dataset.
Such dataset is then used to extract the training data for the ML models and to test them.
The training dataset is generated by sampling the CFD fields, on the finest grid of the reference $11\,\mathrm{m/s}$ case only, along two planes crossing the rotor centre: one parallel and one perpendicular to the terrain, so as to capture the wake-ABL interaction in both directions. This results in a total of $\sim 180k$ cells for the training dataset, as anticipated in Sec.~\ref{sec:NN}. Deliberately limiting the training data to a single grid and a single operating point keeps the high-fidelity cost to the minimum, and turns every coarser grid and every other inflow velocity into an independent test of the closure rather than an additional training case.